% !PRINTED 380
\weberpage{380}
%<<CONT
rational durch diese $a_{k,i}$ bestimmt, und aus (16) und (17) erhält man sodann $\xi_{4i}$, $\xi_{5i}$, $\xi_{6i}$, $\xi_{7i}$, $x_4$, $x_5$, $x_6$, $x_7$.

Durch Substitution der $\xi_1$, $\xi_2$, $\xi_3$ in die Gleichung (4) erhält man die Gleichung einer Curve vierter Ordnung, deren Co\"efficienten rationale Functionen der $a_{k,i}$ sind, und die Discriminante dieser Gleichung kann nicht identisch verschwinden, weil man umgekehrt, wie wir gesehen haben, aus der Gleichung einer Curve vierter Ordnung ohne singulären Punkt ein Gleichungssystem (18), (19), (20) ableiten kann.

Aus den Gleichungen (18), (16) folgen dann auch die Formeln (7), und die daraus durch Vertauschung von 4 mit 5, 6, 7 hervorgehenden, woraus zu schliessen ist, dass $x_4$, $x_5$, $x_6$, $x_7$ zusammen mit $x_1$, $x_2$, $x_3$ ein vollständiges Siebener-System bilden.

\S.\,101.

\gesperrt{Die Galois'sche Gruppe des Doppeltangentenproblems.}

Die Sätze, die wir abgeleitet haben, sind ausreichend, um die \gesperrt{Galois'sche} Gruppe der algebraischen Gleichung zu bestimmen, von der die Doppeltangenten abhängen. Wir betrachten hierbei als \gesperrt{Rationalitätsbereich den Körper, der aus allen rationalen Functionen der 14 Verhältnisse der Co\"efficienten einer allgemeinen ternären Form $4^{\text{ten}}$ Grades} besteht, worin diese Co\"efficienten als \gesperrt{unabhängige Variable} gelten. Die Gleichung $28^{\text{sten}}$ Grades können wir uns dann etwa in der Weise gebildet denken, dass wir als Unbekannte die Abscissen der Schnittpunkte der Doppeltangenten mit einer beliebigen festen geraden Linie $L$ betrachten. Durch die Wurzeln $\xi$ dieser Gleichung, die wir die \gesperrt{Doppeltangentengleichung} nennen, sind dann die Doppeltangenten rational darstellbar.

Benutzt man ein \gesperrt{Cartesisches} Coordinatensystem $x$, $y$, dessen $x$-Axe die Linie $L$ ist, so erhält die Gleichung einer Doppeltangente die Gestalt
\weq{1}{y = \Theta\,(x - \xi),}
und die Doppeltangentengleichung kann gebildet werden, wenn $F(x,y) = 0$ die Gleichung der Curve vierter Ordnung ist, indem man die Bedingungen aufsucht, dass die Function von $x$
%CONT>>
